
\subsubsection{Вычисление нагрузок}

\epigraph{
    Feel what the wave is doing.\\
    Then accept its energy.\\
    Get in sync.\\
    Then charge with it.
}{\textit{-- Бодхи}\\«На гребне волны» (1991)}


После того как волна прошла трансформацию на мелководье (шоалинг, разрушение, рефракция), нас интересует уже не столько сама высота волны, сколько то, какой энергетический поток она несет к берегу и какая его доля будет участвовать в его разрушении. 
Теоретически это описывается через понятие волновой мощности и ее проекции на выбранное направление -- в нашем случае на нормаль к береговой линии.

В линейной теории регулярной поверхностной гравитационной волны средняя полная энергия (сумма кинетическоой и потенциальной) на единицу площади поверхности пропорциональна квадрату амплитуды волны или, в нашем случае, значимой высоты волны. 
Для спектрального описания волнения удобно использовать значимую высоту $H_s$ и пиковый период $T_p$, которые являются интегральными характеристиками энергетического спектра. 
В установившейся волне поток энергии переносится групповой скоростью вдоль направления распространения волны, так что поток энергии на единицу длины фронта (мощность волны) может быть выражен в виде формулы вида~\cite{Muraleedharan2023Spectral}:
\begin{equation}\label{eq:wave_power}
    P \propto \rho\,g^2\,H_s^2\,T_p,
\end{equation}
где $\rho$ -- плотность воды, $g$ -- ускорение свободного падения, а коэффициент пропорциональности определяется конкретной формой спектра.

Одна из часто используемых инженерных форм формулы~\eqref{eq:wave_power} записывается как~\cite{guillou2020estimating}:
\begin{equation}
    P = \frac{\rho\,g^2\,H_s^2\,T_p}{64\pi}.
\end{equation}

Величина $P$ имеет размерность мощности на единицу длины фронта (Вт/м) и характеризует интенсивность волнового потока вдоль направления распространения волнового луча.
В таком определении $P$ относится к направлению распространения волны и показывает, сколько энергии переносится в единицу времени через отрезок фронта длиной 1 м, ориентированный поперек луча. 

Для задач прибрежной динамики этого недостаточно, поскольку береговая линия имеет собственное направление, и только та часть потока, которая ориентирована перпендикулярно к берегу, непосредственно участвует в его разрушении и деформации.
Пусть $\theta$ -- угол между направлением распространения волны и нормалью к береговой линии в прибрежной точке (значение $\theta = 0$ определяет строго перпендикулярный накат на берег). 
Вектор потока энергии $\vec{P}$ направлен вдоль луча, а нас интересует его проекция на нормаль к берегу $\vec{n}$. 
В рамках линейной теории и квазистационарных условий эта проекция задается скалярным произведением~\cite{GalalTakewaka2011}:
\begin{equation}
    P_n = |\vec{P}|\cos\theta = P\cos\theta,
\end{equation}
где $P_n$ -- компонент потока энергии, направленный вдоль $\vec{n}$, то есть на берег. 

В прибрежной энергетике этот показатель часто обозначается как CWEF (Coastal Wave Energy Flux)~\cite{Ghanavati2023_wave_surge_coastlines}:
\begin{equation}\label{eq:CWEF}
    \text{CWEF} = P\cos\theta.
\end{equation}

Таким образом, CWEF -- это мощность волны, пересчитанная с вдольлучевого направления на нормаль к берегу. 
Если волна идет почти параллельно берегу ($\theta \approx 90^\circ$), $\cos\theta \approx 0$, при значительном $H_s$ реальное воздействие на данный профиль по нормали оказывается небольшим. 
Если же волна приходит близко к нормали ($\theta \to 0$), $\cos\theta \to 1$, то практически весь поток энергии работает на деформацию берега и изменение его профиля.

С физической точки зрения величина $\text{CWEF}$ показывает, какая мощность волнового потока в среднем падает на один метр береговой линии в направлении, способствующем разрушению и переотложению наносов~\cite{Pang_2023_CoastalErosionClimate}.

Для заданного $H_s$ изменение угла $\theta$ может радикально изменить $\text{CWEF}$, хотя локальная высота волны у берега остается той же. 
Поэтому с точки зрения прибрежной морфодинамики два состояния с одинаковым $H_s$, но различным $\theta$, неравноценны: при меньшем $\theta$ (более перпендикулярный подход) CWEF выше, и берег подвергается большей энергетической нагрузке.

Если рассматривать CWEF как функцию времени $\text{CWEF}(t)$, то интеграл
\begin{equation}
    E_\text{int} = \int_{t_1}^{t_2} \text{CWEF}(t)\,dt
\end{equation}
имеет размерность энергии на единицу длины берега (Дж/м) и может интерпретироваться как суммарная энергия, доставленная волнами к данному участку берега за период $[t_1, t_2]$. Нормируя эту величину на длительный интервал, можно получать оценки «волновой экспозиции» участка берега, которые далее могут сопоставляться с наблюдаемой интенсивностью эрозионных процессов и использоваться в индексах прибрежной уязвимости и других агрегированных показателях.
